% =====================================================================
% artigo.tex — Template de artigo acadêmico ABNT/Unifor PPGD
% Compilação: latexmk -pdf -xelatex artigo.tex
%        ou: ./build.sh
% =====================================================================
\documentclass[
    article,        % modo artigo (não livro/tese)
    12pt,
    a4paper,
    oneside,        % anverso apenas
    brazil,
    sumario=tradicional
]{abntex2}

% =====================================================================
% preamble.tex — Configuração ABNT estrita para PPGD/Unifor
% Doutorado em Direito Constitucional — Kalyl Lamarck
% Compatível com xelatex (suporte UTF-8 + fontes nativas)
% =====================================================================

% --- Encoding e idioma ---------------------------------------------------
% xelatex le UTF-8 nativo (inputenc nao necessario)
% babel ja carregado pela classe abntex2 (option 'brazil')

% --- Fonte Times New Roman (msttcorefonts) -------------------------------
\usepackage{fontspec}
\setmainfont{Times New Roman}       % corpo
\setsansfont{Times New Roman}       % forca serifada tambem em sans
\setmonofont{Liberation Mono}        % monoespacada (em codigo)

% --- Tamanho 12pt em tudo (incluindo headings) ---------------------------
% abntex2 usa, por default, \large/\Large nos headings. Forçamos \normalsize
% para conformidade com NBR 14724 (corpo e títulos em 12pt).
\renewcommand{\ABNTEXchapterfont}{\bfseries}
\renewcommand{\ABNTEXchapterfontsize}{\normalsize}
\renewcommand{\ABNTEXsectionfont}{\bfseries}
\renewcommand{\ABNTEXsectionfontsize}{\normalsize}
\renewcommand{\ABNTEXsubsectionfont}{\bfseries}
\renewcommand{\ABNTEXsubsectionfontsize}{\normalsize}
\renewcommand{\ABNTEXsubsubsectionfont}{\bfseries}
\renewcommand{\ABNTEXsubsubsectionfontsize}{\normalsize}

% Forçar CAIXA ALTA em chapter (seção primária) — texto digitado ficará upper
% via comando \MakeUppercase no título; abntex2 já força isso por default.

% --- Espaçamento entre linhas: 1,5 (corpo); simples (notas/refs) ---------
\OnehalfSpacing

% --- Margens ABNT 3/2/3/2 cm ---------------------------------------------
\usepackage[
    a4paper,
    top=3cm,
    bottom=2cm,
    left=3cm,
    right=2cm,
    bindingoffset=0pt
]{geometry}

% --- Recuo de primeira linha: 1,5 cm -------------------------------------
\setlength{\parindent}{1.5cm}
\setlength{\parskip}{0pt}

% --- Citações: estilo autor-data ABNT (abntex2cite) ----------------------
% [alf] = autor-data ('alfabético' no jargão ABNT)
% [abnt-emphasize=bf] = título da obra em negrito
% [abnt-etal-list=0] = sem 'et al.' (lista todos os autores)
% [abnt-and-type=&] = separador entre autores no texto
\usepackage[
    alf,
    abnt-emphasize=bf,
    abnt-etal-list=0,
    abnt-and-type=&,
    abnt-thesis-year=both,
    abnt-repeated-author-omit=yes
]{abntex2cite}

% --- Hyperref sem cores (impressão acadêmica preto-e-branco) -------------
\usepackage{hyperref}
\hypersetup{
    colorlinks=false,
    pdfborder={0 0 0},
    linkcolor=black,
    citecolor=black,
    urlcolor=black,
    pdftitle={\@title},
    pdfauthor={\@author},
    pdflang=pt-BR
}

% --- Tabelas e figuras ---------------------------------------------------
\usepackage{graphicx}
\usepackage{booktabs}
\usepackage{tabularx}
\usepackage{caption}

% Legendas em TNR 12pt, alinhadas à esquerda
\captionsetup{
    font={normalsize},
    labelfont={bf},
    justification=raggedright,
    singlelinecheck=false
}

% --- Notas de rodapé: TNR 10pt, espaçamento simples ----------------------
\renewcommand{\footnotesize}{\fontsize{10pt}{12pt}\selectfont}
\setlength{\footnotesep}{14pt}        % espaço entre notas
\renewcommand{\footnoterule}{\rule{4cm}{0.4pt}}

% --- Outras configurações estéticas --------------------------------------
\setlength{\afterchapskip}{1.5\baselineskip}
\setlength{\beforechapskip}{0pt}

% --- Renomear titulo da bibliografia para REFERENCIAS (NBR 6023) ---------
% Em modo 'article' nao ha \chapter, entao redefine \bibsection
% para usar \section* em vez do \chapter* default do abntex2.
\renewcommand{\bibname}{REFERÊNCIAS}
\renewcommand{\bibsection}{%
    \section*{\bibname}%
    \addcontentsline{toc}{section}{\bibname}%
}

% --- Variáveis institucionais (usar nos textos) --------------------------
\newcommand{\universidade}{Universidade de Fortaleza}
\newcommand{\programa}{Programa de Pós-Graduação em Direito Constitucional}
\newcommand{\nivel}{Doutorado em Direito Constitucional}
\newcommand{\autorprincipal}{Kalyl Lamarck Filgueiras Estanislau Costa}


% --- Metadados do artigo -------------------------------------------------
\title{Título do artigo em português}
\author{
    \textbf{Nome do Autor 1}\thanks{Filiação institucional do autor 1, e-mail.}
    \and
    \textbf{Nome do Autor 2}\thanks{Filiação institucional do autor 2, e-mail.}
}
\date{}    % suprimir data automática

% =====================================================================
\begin{document}
\selectlanguage{brazilian}
\frenchspacing

% --- Título centralizado, TNR 12pt negrito ---
\begin{center}
\textbf{\thetitle}
\end{center}

\vspace{1ex}

\begin{center}
\textit{\textbf{Título do artigo em inglês}}
\end{center}

\vspace{1ex}

% --- Autores alinhados à direita, CAIXA ALTA negrito ---
\begin{flushright}
\MakeUppercase{\textbf{Nome do Autor 1}}\footnote{Filiação institucional, e-mail, ORCID.}\\
\MakeUppercase{\textbf{Nome do Autor 2}}\footnote{Filiação institucional, e-mail, ORCID.}
\end{flushright}

\vspace{2ex}

% --- RESUMO -------------------------------------------------------------
\noindent\textbf{RESUMO}

\noindent
Texto do resumo em parágrafo único, máximo 150 palavras, justificado, sem recuo, espaçamento simples. Apresenta tema, objetivo, método, resultados e conclusão.

\vspace{1ex}

\noindent\textbf{PALAVRAS-CHAVE}

\noindent
Termo 1; termo 2; termo 3; termo 4; termo 5.

\vspace{1ex}

% --- ABSTRACT -----------------------------------------------------------
\noindent\textbf{ABSTRACT}

\noindent
Abstract text in single paragraph, max 150 words, justified, no indent, single spacing. Presents topic, objective, method, results, and conclusion.

\vspace{1ex}

\noindent\textbf{KEYWORDS}

\noindent
Term 1; term 2; term 3; term 4; term 5.

\vspace{2ex}

% --- INTRODUÇÃO ---------------------------------------------------------
% Em modo article, usa-se \section* para INTRODUÇÃO/CONCLUSÃO/REFERÊNCIAS
% (sem indicativo numérico) e \section para "1 ...", "2 ..."
\section*{INTRODUÇÃO}
\addcontentsline{toc}{section}{INTRODUÇÃO}

% INTRODUÇÃO — sem citações, conforme NBR 14724

Os sertões e as caatingas do Ceará e do Rio Grande do Norte abrigam, entre 2004 e 2025, uma das mais extensas expansões de parques eólicos do continente americano. A arquitetura institucional reproduz, sob forma renovável, a racionalidade econômica que orienta a alocação geopolítica de riscos.

Pergunta-se se a expansão eólica no Nordeste brasileiro configura transição energética genuína nos termos da racionalidade ambiental, ou se reproduz, sob discurso verde, lógica análoga à formalizada no início dos anos 1990 sobre territórios de baixa renda.

O objetivo deste artigo consiste em demonstrar que o arranjo fiscal, contratual, territorial, humano e multiespécie do setor eólico cearense e potiguar constitui padrão de captura verde, composto por cinco vetores simultâneos.

Justifica-se a investigação por duas dimensões convergentes. No plano social, comunidades quilombolas, assentadas e pesqueiras têm sido afetadas sem consulta livre, prévia e informada. No plano acadêmico, a literatura jurídica brasileira examina cada vetor de maneira fragmentária, sem amarrar a linearidade institucional que liga a doutrina inicial à expansão eólica nordestina contemporânea.

Metodologicamente, a investigação adota epistemologia dogmática jurídica de tradição analítica, métodos hermenêutico-argumentativo e dedutivo, e técnicas de pesquisa bibliográfica, documental e jurisprudencial. As fontes são primárias e secundárias; terciárias proibidas.

O artigo organiza-se em dois capítulos. O Capítulo 1 constrói a régua conceitual-normativa da transição energética. O Capítulo 2 testa empiricamente essa régua sobre a eólica do Ceará e do Rio Grande do Norte.


% --- SEÇÕES NUMERADAS ---------------------------------------------------
% Capítulo 1 — Régua conceitual-normativa

\section{\MakeUppercase{Transição energética e seus critérios de ruptura}}

\subsection{Racionalidade ambiental como critério epistêmico}

Conforme \citeonline{leff2006}, a racionalidade ambiental constitui-se como contracorrente da racionalidade econômica dominante e opõe, à redução da natureza a capital fungível, a reapropriação social do mundo natural por meio dos saberes locais.

A incomensurabilidade é o mecanismo central que \citeauthoronline{leff2006} identifica para distinguir as duas racionalidades. Pois constata-se que os valores simbólicos, ecológicos, epistemológicos e políticos da natureza e da cultura resistem à homologação ao equivalente universal do mercado.

\subsection{Decrescimento e bem viver}

Sob outro ângulo, \citeauthoronline{latouche2009} sustenta que o decrescimento não constitui antimodernismo, mas recusa do imaginário econômico segundo o qual o crescimento se apresenta como única forma social aceitável.

Em chave complementar, \citeauthoronline{acosta2016} opera no plano relacional, propondo o bem viver como suficiência relacional fixada pela comunidade em sua relação com o território.

% Capítulo 2 — Aplicação empírica

\section{\MakeUppercase{A eólica no Nordeste como captura da transição energética}}

\subsection{Vetor fiscal e acumulação por despossessão}

Aplicada ao recorte cearense e potiguar, a régua de \citeauthoronline{harvey2010} opera com precisão sobre a engrenagem que articula \textit{Regime Especial de Incentivos para o Desenvolvimento da Infraestrutura} (REIDI) à isenção SUDENE.

Com efeito, \citeauthoronline{piketty2014} demonstra que a desigualdade fundamental em que a taxa de retorno do capital supera a taxa de crescimento da economia constitui condição estrutural do capitalismo avançado.

\subsection{Vetor contratual e expulsão pelos contratos eólicos}

Sob esse prisma, \citeauthoronline{ostrom1990} formula sete princípios construtivos para a sustentação de instituições comunais de longa duração. Aplicada ao recorte, a régua de Ostrom diagnostica que os contratos eólicos extinguem o sétimo princípio, ao remover a possibilidade de escolha coletiva entre comunidade anfitriã e operadora.


% --- CONCLUSÃO ----------------------------------------------------------
\section*{CONCLUSÃO}
\addcontentsline{toc}{section}{CONCLUSÃO}

% CONCLUSÃO — sem citações novas, conforme NBR 14724

Retomado o objetivo proposto na Introdução, sintetiza-se o rendimento do teste empírico dos cinco vetores de captura. A investigação reuniu, no Capítulo 1, cinco réguas conceituais-normativas que operam em ângulos distintos sobre a transição energética.

No Capítulo 2, essa régua composta foi aplicada à expansão da energia eólica no Ceará e no Rio Grande do Norte entre 2004 e 2025. Demonstra-se a linearidade institucional que liga a doutrina formalizada na década de 1990 ao arranjo eólico contemporâneo.

Sustenta-se, em primeiro achado, que a linearidade institucional opera factualmente, e não por analogia retórica. Em segundo achado, os cinco vetores operam simultaneamente.

Outrossim, em terceiro achado, o discurso de transição verde, quando desprovido de critério de racionalidade ambiental substantiva, configura mecanismo de redistribuição fiscal sem participação política e sem proteção multiespécie.

Propõe-se, como fechamento, que uma transição energética compatível com a racionalidade ambiental exige consulta livre, prévia e informada obrigatória, redistribuição fiscal ancorada em indicadores sociais municipais e monitoramento público independente.


% --- REFERÊNCIAS --------------------------------------------------------
\section*{REFERÊNCIAS}
\addcontentsline{toc}{section}{REFERÊNCIAS}
\renewcommand{\bibsection}{}
\bibliography{referencias/referencias_doutorado}

\end{document}
